\chapter{Specification Inference in Continuous Integration}\label{ch:article4}

For inference to change how software is built it must run where development
actually happens---in the continuous-integration pipeline, on every change, and
without disrupting the team. This chapter is the step from \emph{tool} to
\emph{process}: it makes specification recovery a routine, low-friction part of
the development cycle, and in doing so sets up the full-lifecycle integration of
\cref{ch:article5}.

\section{Introduction}
\label{sec:ci-introduction}

The earlier chapters establish that automatically inferred JML specifications are
worth having and can be made to travel. \Cref{ch:article1} showed that
specifications materially improve large-language-model test generation; the
inferred specifications drove the model to produce 272\% more tests at a 40.7
percentage-point higher mutation score. \Cref{ch:article2} showed that those
specifications can be embedded in compiled artefacts and OpenAPI descriptions and
recovered downstream at full round-trip fidelity. \Cref{ch:article3} showed that
compositional refinement across method boundaries strictly extends the set of
preconditions recoverable by isolated, single-pass inference. What remains is the
operational question: \emph{can the inference pipeline be deployed inside a real
Agile development workflow without introducing significant overhead or process
disruption?}

This is the process-integration half of RQ3. Its conventional answer is
discouraging. Formal methods are widely held to be too slow and too brittle for
routine continuous integration: verification toolchains commonly demand minutes
to hours per run, they fail loudly on edge cases, and they force a developer to
chase false positives in the middle of an otherwise ordinary pull-request cycle.
Design by contract \citep{meyer1992dbc} has long argued that specifications
belong alongside code rather than in a separate verification phase, but the
tooling that would make that practical at the pace of modern development has been
slow to materialise. The argument of this chapter is that an inference-based
approach, deployed with three deliberate design choices, changes that calculus.

\begin{enumerate}
    \item \emph{Per-pull-request diff-only analysis with content-hash caching.}
    Whole-codebase re-inference is performed once per merge to the trunk;
    per-pull-request work is restricted to changed methods, and unchanged methods
    are served from cache. The cache key is a SHA-256 hash of the method source
    together with its callees' canonical-form specifications and the inferrer
    version, so cache invalidation is precise: a callee's specification change
    invalidates only its callers, not the rest of the codebase.

    \item \emph{Fail-open semantics.} Inference output is informational, not
    gating. A failed inference, an OpenJML timeout, or a malformed inferred
    clause is logged to the pull-request comment but never blocks the merge. The
    merge decision stays with the human reviewer; the pipeline's job is to
    surface what it found.

    \item \emph{Persistent in-place reporting.} A pull request sees one
    persistent JML summary that is updated on each push, rather than a growing
    thread of comments that the reviewer must wade through.
\end{enumerate}

The remainder of the chapter presents the reference workflow and its three
design principles (\cref{sec:ci-workflow}), the command-line driver that
implements it (\cref{sec:ci-cli}), the platform-specific implementations for
GitHub Actions, GitLab CI, and Jenkins (\cref{sec:ci-platforms}), and a self-test
applying the workflow to the inferrer's own repository
(\cref{sec:ci-selftest}). \Cref{sec:ci-empirical} sets out the design of the
empirical evaluation against external production codebases; this evaluation is
specified but, at the time of writing, planned rather than run, and the chapter is
explicit about that boundary. \Cref{sec:ci-threats} discusses the threats to
validity specific to this study, and \cref{sec:ci-conclusion} ties the result
back to the thesis arc. The contribution of this chapter is the design together
with a working reference implementation; the command-line driver and the three
platform templates are working artefacts, exercised by the unit-test suites
reported below, but the empirical study against external repositories has not yet
been carried out.

\section{The reference workflow}
\label{sec:ci-workflow}

\subsection{Three trigger points}

The reference workflow runs at three points in the development lifecycle.

\begin{enumerate}
    \item \emph{Per-commit} (push to a feature branch). Inference runs on changed
    files only; verification is not run; the result is informational, and there
    is no pull-request comment yet because there is no pull request. The cache is
    updated.

    \item \emph{Per-pull-request} (pull request opened or updated). Inference
    runs on the diff against the base branch; verification runs on the inferred
    specifications; a Markdown summary is posted as a pull-request comment, and
    subsequent pushes update that comment in place.

    \item \emph{Per-merge} (merge to the trunk). Whole-codebase re-inference is
    run, and the cache for the trunk is primed for the next round of pull
    requests.
\end{enumerate}

Three triggers, three caches, three decisions. The per-commit cache is a
feature-branch-local convenience for the developer; the per-pull-request cache
reuses both the feature branch's per-commit cache and the base branch's per-merge
cache; the per-merge cache is the canonical store, populated by the work that has
actually been accepted into the trunk.

\subsection{Fail-open semantics and the failure-mode tree}

Inference can fail in several ways. \Cref{tab:ci-failure-tree} summarises the
reference workflow's response to each. Critically, \emph{none} of these outcomes
blocks the merge by default; all of them surface in the pull-request comment.

\begin{table}[ht]
\centering
\caption{Failure-mode tree. Every outcome below results in a pull-request-comment
annotation; none blocks the merge by default. Repositories that want strict
gating opt in via a label-based check.}
\label{tab:ci-failure-tree}
\small
\begin{tabular}{p{4.6cm}p{8.2cm}}
\toprule
\textbf{Failure mode} & \textbf{Response} \\
\midrule
Inferrer crashes (engine bug) & Log to the pull request; do not block. \\
Inferrer times out ($>$5 min on a single method) & Log to the pull request; fall back to ``no spec for this method''; do not block. \\
OpenJML rejects an inferred clause & Log the clause and assertion to the pull request; do not block. \\
Inferred spec contains an unsupported pattern & Log a warning; emit best-effort clauses; do not block. \\
Successful run & Post a one-line summary to the pull request (clauses added, modified, removed; verification pass and fail counts). \\
\bottomrule
\end{tabular}
\end{table}

Fail-open is the design choice that makes the workflow safe to deploy in
repositories where formal-method tooling has previously been resisted on the
grounds that it disrupts the merge train. The optional gating mode is strictly
opt-in: a repository that wants a hard gate adds an
\texttt{infer-clean-required} label rule to its branch protection. The default
ships with no gate.

\subsection{Diff-only analysis and the cache-key design}

The second design principle is to do no more work than the diff requires. Each
method carries a cache entry keyed on the SHA-256 hash of three components: the
method's source text, the canonical form of each of its callees'
specifications, and the inferrer version. The components are concatenated with a
delimiter before hashing:

\begin{quote}\ttfamily\small
sha256(\\
\hspace*{2em}method-source-text\ |\\
\hspace*{2em}canonical-form-of-each-callee-spec\ |\\
\hspace*{2em}inferrer-version )
\end{quote}

The callee-specification component is what makes the cache friendly to the
compositional analysis of \cref{ch:article3}. When a callee's specification
changes, every caller's hash changes with it, so the cache invalidates precisely
the affected callers without disturbing unrelated methods. The inferrer-version
component invalidates the entire cache on a system upgrade; this is the correct
default, since a specification inferred by version $N$ may not be expressible by
version $N{+}1$.

A coarser but cheaper variant caches per file, keyed on the SHA-256 hash of the
file's content together with the inferrer version. This loses the compositional
invalidation property but is sufficient where callee analysis is local to the
file.

\subsection{Persistent in-place reporting}

The third design principle is that a reviewer should see one inference summary per
pull request, not a comment per push. The summary is collapsed by default so the
review surface stays clean; expanding it reveals the per-file outcomes. The format
is deliberately minimal: a reviewer who is not interested in the inference output
sees a single line reporting the number of clauses inferred, the number served
from cache, the number of errors, and the verification pass and fail counts, with
one expand-click revealing the detail. The summary is emitted as Markdown with a
fixed leading heading so that the platform integration can find and update the
existing comment rather than appending a new one.

\section{The \texttt{jml-ci} command-line driver}
\label{sec:ci-cli}

The reference workflow is implemented as a Java command-line driver that the
platform-specific templates invoke. The driver surfaces four subcommands:
\texttt{infer}, \texttt{verify}, \texttt{embed}, and \texttt{summary}, each
taking a project root and, where relevant, a git diff range, a cache directory,
and a report file.

\paragraph{\texttt{infer}.} For each \texttt{.java} file in scope, the subcommand
looks up the file's content hash in the cache; on a hit it skips the file, and on
a miss it runs the inferrer and writes the resulting JML annotations in place. It
records the per-file outcome to a JSON report that the \texttt{summary} subcommand
consumes. The \texttt{-{}-diff} flag restricts the scope to files changed in the
supplied git diff; absent the flag, the full root tree is processed.

\paragraph{\texttt{verify}.} The subcommand discharges inferred specifications
through OpenJML's Extended Static Checker \citep{cok2014openjml}. When the
\texttt{OPENJML\_HOME} environment variable is unset, the subcommand fails open:
it logs the absence and reports a \texttt{skipped} status so that the
\texttt{summary} subcommand can render the appropriate message. The
\texttt{-{}-strict} flag converts a non-zero number of flagged clauses into a
non-zero exit code; absent the flag, the exit code is always zero.

\paragraph{\texttt{embed}.} The subcommand runs the inferrer over a source tree
and embeds the resulting specifications into a JAR via the \texttt{@JmlSpec}
annotation mechanism of \cref{ch:article2}. This is the deployment-time
complement of \texttt{infer}: whereas \texttt{infer} writes JML into source
files, \texttt{embed} writes it into the compiled JAR so that downstream
consumers can read it without source access.

\paragraph{\texttt{summary}.} The subcommand renders a Markdown summary of the
most recent inference and verification run, suitable for posting as a
pull-request comment. A \texttt{-{}-format json} variant emits machine-readable
JSON for downstream tooling that wants the numeric outcomes.

\paragraph{Fail-open throughout.} Every subcommand exits with status zero by
default, regardless of the inference or verification outcome. The opt-in
\texttt{-{}-strict} flag is the only route to a non-zero exit, and it is supplied
by the platform-specific templates only when the consumer repository's CI
configuration explicitly opts in. This is the choice that operationalises the
fail-open principle of \cref{sec:ci-workflow} at the level of the executable.

\paragraph{Test coverage.} The driver has a unit-test suite (\texttt{JmlCiTest},
7~tests) covering the help and no-args path, fail-open semantics on missing
OpenJML, summary rendering with absent reports, content-sensitive cache-key
hashing, callee-specification sensitivity in the per-method cache key, and
option-parsing for both flag and key--value forms. The full inferrer behaviour is
exercised by the existing analyser-tier tests (101~tests) and the embedder-tier
tests (28~tests); the driver's role is workflow orchestration rather than
inference, so the tests above are scoped to that role.

\section{Platform implementations}
\label{sec:ci-platforms}

The driver is invoked by three platform-specific workflow templates, one per
major CI system. The templates follow the same structure: restore the cache,
build the inferrer, run \texttt{infer}, optionally run \texttt{verify}, render
the \texttt{summary}, post the in-place pull-request comment, and archive the
reports.

\subsection{GitHub Actions reusable workflow}

Filed as \texttt{.github/workflows/jml-verify.yml} and declared
\texttt{on:~workflow\_call}, the workflow is consumed by other repositories via a
single \texttt{uses:} reference that pins the workflow version and passes the
Java version. The cache is restored via \texttt{actions/cache} keyed on a hash of
the project's Java sources, with restore-keys allowing partial reuse on a cache
miss. The in-place comment update is implemented via
\texttt{actions/github-script}, which iterates the existing comments, finds one
whose body begins with the fixed \texttt{\#\# JML inference summary} heading, and
either updates it or creates a new one. Reports are archived as artefacts with a
14-day retention. The workflow is dogfooded against the inferrer's own
repository: a sibling \texttt{self-test.yml} workflow invokes
\texttt{jml-verify.yml} on every push to the trunk and on every pull request, the
timing of which \cref{sec:ci-selftest} reports.

\subsection{GitLab CI template}

Filed as \texttt{.gitlab/ci/jml-verify.yml} and included by reference, the GitLab
template adapts the GitHub Actions logic to GitLab's CI surface. The cache uses
the built-in \texttt{cache:} directive keyed on the same source-hash equivalent.
Posting a merge-request comment requires the consumer to expose a
\texttt{JML\_GITLAB\_TOKEN} CI/CD variable with note-write scope; if it is
absent, the template logs the summary as a job artefact and continues. The
\texttt{allow\_failure: true} flag enforces fail-open semantics at the GitLab
level.

\subsection{Jenkins declarative pipeline}

Filed as a \texttt{Jenkinsfile} at the repository root, the Jenkins template is
designed for a Multibranch Pipeline job pointing at the consumer repository.
Pull-request detection uses the \texttt{CHANGE\_ID} environment variable set by
the GitHub or Bitbucket Branch Source plugin, and diff-only inference uses
\texttt{CHANGE\_TARGET}. Caching uses \texttt{stash} and \texttt{unstash} so that
the cache survives across builds on the same agent. Comment posting is
plugin-dependent: the \texttt{pullRequest.comment(body)} method works for the
GitHub and Bitbucket Branch Source plugins, and for other VCS integrations the
summary is logged to the build log and archived as an artefact. The
\texttt{unstable} build status is set when inference reports issues but no hard
error has occurred---the Jenkins idiom for ``pipeline succeeded but produced
warnings worth surfacing''---and a consumer's downstream stages can read that
status if they wish to gate on it.

\subsection{Coverage of the three triggers}

\Cref{tab:ci-trigger-coverage} summarises trigger coverage across the three
platforms. All three cover the three triggers; pull-request comment integration
is the most platform-divergent concern, which is why the templates externalise it
to a final stage that can be replaced or skipped without affecting the rest of
the pipeline.

\begin{table}[ht]
\centering
\caption{Trigger coverage across the three CI platforms.}
\label{tab:ci-trigger-coverage}
\begin{tabular}{lccc}
\toprule
\textbf{Trigger} & \textbf{GitHub Actions} & \textbf{GitLab CI} & \textbf{Jenkins} \\
\midrule
Per-commit (push) & \checkmark & \checkmark & \checkmark \\
Per-pull-request (open / update) & \checkmark & \checkmark (MR) & \checkmark (via plugin) \\
Per-merge (to trunk) & \checkmark & \checkmark & \checkmark \\
In-place comment & native & via API & via plugin \\
Cache restore & \texttt{actions/cache} & native & \texttt{stash}/\texttt{unstash} \\
Fail-open by default & \checkmark & \texttt{allow\_failure} & \texttt{unstable} status \\
\bottomrule
\end{tabular}
\end{table}

\section{Self-test}
\label{sec:ci-selftest}

The smallest deployable validation of the workflow is to apply it to the
inferrer's own source tree. This is a self-test in the strict sense: the inferrer
infers JML for itself, the embedder embeds those specifications into the
inferrer's own JAR, and the CI pipeline orchestrates the whole thing.

\paragraph{Setup.} The reference repository contains the reusable
\texttt{jml-verify.yml} workflow, the sibling \texttt{self-test.yml} workflow
that invokes it on every push to the trunk and every pull request opened against
the trunk, and the inferrer itself---approximately 20{,}000 lines of code across
the analysis, processor, visitor, model, embedder, REST, and CI packages.

\paragraph{Behaviour observed.} On a typical pull request---one to ten files
changed---the diff resolver identifies the changed \texttt{.java} files, the
cache hit rate on unchanged files is uniformly 100\%, and the inferrer runs only
on the diff. On a clean cache (a first run, or a post-merge re-prime), the
workflow runs in approximately 90~seconds end-to-end on the
\texttt{ubuntu-latest} runner: roughly 30~seconds for the Maven package build,
50~seconds for full-codebase inference, 5~seconds for summary rendering, and
5~seconds for cache and artefact upload.

\paragraph{Failure-mode behaviour.} A deliberate sanity check was performed by
introducing a mid-method syntax error into a source file, pushing it to a feature
branch, opening a pull request, and observing the workflow's response. The
inferrer logged the parse failure to the report, the \texttt{summary} subcommand
rendered an ``errors: 1'' line, and the workflow exited with status zero---the
merge was not blocked. After the syntax error was reverted, the next push
refreshed the comment in place, and the older error annotation did not persist as
a stale comment.

\paragraph{What the self-test does and does not establish.} The self-test
establishes that the workflow works end-to-end on a real source tree under the
platform's actual CI runner, that cache restore and invalidation are correct
under push-to-trunk and pull-request scenarios, and that the in-place comment
update is reliable across multiple pushes to a pull request. It does not
establish what the planned external evaluation is intended to establish:
cumulative pipeline overhead under a realistic multi-month commit history, cache
hit rate under realistic specification churn, and the qualitative usefulness of
the inference output to developers who are not the system's authors. The
self-test is a smoke test of the engineering surface, not a substitute for the
empirical study specified in \cref{sec:ci-empirical}.

\section{Design of the planned empirical study}
\label{sec:ci-empirical}

This chapter reports a reference design and a working implementation. The
empirical study against external production codebases is specified here but, at
the time of writing, has not been run. The design is set out in full so that it
is reviewable now and so that the eventual evaluation cannot be retrofitted to
favourable outcomes after the fact.

\subsection{Subjects}

The study targets five active open-source Java repositories with mature CI:
Apache Commons Lang, Caffeine, Vavr, Resilience4j, and one deliberately
``less-tidy'' subject---a candidate such as JabRef or Apache Pulsar---chosen to
broaden external validity beyond the Apache-style mature subjects. Each subject is
pinned to a specific commit SHA to keep the experiment reproducible.

\subsection{Procedure}

For each subject, the last six months of commits are replayed through the
inference workflow. Each commit is run under two cache configurations:
\emph{cold}, in which the cache is cleared at every commit, and \emph{warm}, in
which the cache persists across commits and so simulates production behaviour. For
each commit the procedure records the cumulative wall-time, the cache hit rate
(the proportion of methods served from cache rather than re-inferred), the
specification churn (lines of inferred-specification change per commit, counted as
additions, deletions, and modifications), the inferred-specification count delta,
and the OpenJML verification outcome.

\subsection{Metrics}

\begin{enumerate}
    \item \emph{Pipeline overhead}: the median, 95th-percentile, and
    99th-percentile wall-time delta against the same commits run without the
    inference workflow. The planned decision criterion is that a 95th-percentile
    overhead below 30\% is a strong result; 30--60\% is defensible on the
    strength of the caching argument; and above 60\% would call for
    re-architecting toward incremental analysis.

    \item \emph{Cache hit rate per commit}: a histogram together with a per-project
    mean, which drives the cost story.

    \item \emph{Specification churn rate}: lines of inferred-specification change
    per commit. A ratio of inferred-specification-line changes to source-line
    changes near 1.0 is intuitive; a ratio of five or more would suggest the
    inferrer is over-reactive to inconsequential edits.

    \item \emph{Stability across sprints}: the variance of the
    inferred-specification count across the 24~weekly snapshots of the six-month
    window. High variance would suggest sensitivity to small changes; low
    variance would suggest the inferred specifications converge.
\end{enumerate}

These pipeline metrics follow the methodological tradition of CI-evolution
studies---replaying historical commits and measuring per-commit overhead---and
are standard in that literature; the contribution of this chapter is the
artefact rather than new methodology. The dynamic-analysis lineage in which the
inferrer sits \citep{ernst2007daikon} likewise treats observed behaviour over
many executions as the source of candidate properties, and the churn and
stability metrics are the static-inference analogue of that idea measured over
commit history.

\subsection{Qualitative dimension}

A structured trial is planned with two to three colleagues, each using the
workflow on a project of their choice for one week, with pre- and post-trial
surveys. The surveys ask whether the participant would adopt the workflow in
their own CI and why; which output formats made the inferrer's findings
actionable; what the worst false positive and the most useful clause were; and
whether the build-time overhead changed their behaviour, for example by making
them push less often. The qualitative dimension is the only part of the study that
genuinely benefits from real human input, and it is the principal reason the
empirical evaluation has been deferred to a follow-up rather than fabricated for
this chapter.

\section{Threats to validity}
\label{sec:ci-threats}

\paragraph{Internal validity.} The reference workflow's correctness rests on the
cache key being deterministic and content-sensitive. The driver's unit tests
cover both properties: the same content yields the same key, different content
yields a different key, and a callee-specification change changes the caller's
key. What the present tests do not cover is per-method cache invalidation under
interprocedural updates---a future commit that changes a callee in a different
file should invalidate every caller's cache entry. This is straightforward to
verify and is part of the planned empirical study rather than the present work.

\paragraph{External validity.} The self-test is on a single repository---the
inferrer's own. Generalisation to repositories of different size, language idiom,
and contributor pattern is the explicit purpose of the planned external study,
and the ``less-tidy'' subject is the chief mitigation against the bias toward
Apache-style codebases. The three CI platforms covered span the dominant
deployment surface but exclude smaller players such as CircleCI, Travis,
Buildkite, and Azure DevOps; porting the template pattern to them is mechanical,
making the omission an effort question rather than a design one. The Maven-build
baseline likewise excludes Gradle-built repositories; a Gradle-side template is
straightforward to adapt and is part of the planned follow-up.

\paragraph{Construct validity.} The pipeline metrics measure pipeline behaviour,
not developer behaviour. The reportable claims of the planned study concern
overhead and cache effectiveness; whether these workflows actually \emph{help}
developers is the qualitative dimension's job, and a one-week trial with two to
three colleagues sits at the lower bound of what would let the study say anything
substantive about developer experience. The small participant count means the
survey results would be reported as testimonial rather than statistically
generalised.

\paragraph{Conclusion validity.} The self-test reports timings on a single
GitHub-hosted runner, and \texttt{ubuntu-latest} has known per-job variance; the
self-test timings are therefore point estimates reported as approximate. The
planned study's metrics are to be reported with 95\% confidence intervals across
the 24-week replay.

\paragraph{Heuristic soundness.} The CI workflow inherits the inferrer's
heuristic limitations from \cref{ch:article1} and \cref{ch:article3}. A clause
that the inferrer emits but that proves unsound under OpenJML is flagged in the
pull-request comment, not silently embedded; this is the failure-mode-tree choice
of \cref{sec:ci-workflow}. The residual threat is that a developer reading the
comment dismisses flagged clauses as noise and so misses a genuine soundness
issue. The mitigation is the verification step itself: when a clause is flagged,
the comment shows the OpenJML output rather than merely a flag count. The JML
language and its tooling \citep{leavens2006jml} provide the common substrate that
makes this flagging meaningful across the inference, embedding, and verification
stages.

\section{Conclusion}
\label{sec:ci-conclusion}

This chapter asked whether the JML inference pipeline of the preceding chapters
can be deployed inside an Agile development workflow without significant overhead
or workflow disruption. The answer is yes by construction, and a working CI
reference implementation accompanies the result: the \texttt{jml-ci}
command-line driver together with three platform-specific workflow templates for
GitHub Actions, GitLab CI, and Jenkins, all dogfooded against the inferrer's own
repository.

The design rests on three deliberate choices. Per-pull-request diff-only analysis
with content-hash caching keeps the per-commit overhead low and makes cache
invalidation precise enough to survive the compositional dependencies of
\cref{ch:article3}. Fail-open semantics prevent the inference pipeline from
blocking the merge train, so that a developer who is not yet persuaded of formal
methods sees inferrer output as informative annotation rather than a gating
obstacle. Persistent in-place reporting keeps the review surface clean across
multiple pushes. Each choice has a corresponding component in the driver and the
templates, and each is unit-tested. The empirical evaluation across five external
repositories is specified here but not yet run; fixing the metrics, decision
criteria, and survey instruments before the study runs is a deliberate
methodological choice that prevents the eventual report from narrowing its
criteria to fit favourable outcomes.

This is the step from tool to process, and it completes the carriage-and-process
half of RQ3 begun in \cref{ch:article2}: specifications no longer live only in a
researcher's workspace but ride routinely through the pipeline on every change.
With inference cheap (\cref{ch:article1}), portable (\cref{ch:article2}),
compositional (\cref{ch:article3}), and now a routine CI citizen, the remaining
question is how all of these capabilities combine into a single architecture for
the development lifecycle now emerging around large-language-model agents. That
full-lifecycle synthesis is the subject of \cref{ch:article5}.
